\section{方程式と不等式}\label{節:微分方程式不等式}
$x=\cos x$ のような方程式は,解を式で表すことができない.それでも,解が存在することや,解の個数なら調べられるのではないだろうか.また,不等式 $\log x\le x-1$ は,両辺のグラフを描いて比べただけでは証明にならない.

どちらも,左右の差を1つの関数にしてやればよい.方程式 $f(x)=g(x)$ と不等式 $f(x)\ge g(x)$ は,
\[
 h(x)=f(x)-g(x)
\]
と置くと,それぞれ
\[
 h(x)=0,\qquad h(x)\ge0
\]
と書き直せる.どちらも,1つの関数 $h$ の値や符号を調べる問題に帰着するのである.

方程式では,解の存在と個数を分けて示す.\,$h$ が区間で連続で,区間の両端で $h$ の符号が異なれば,中間値の定理により $h(x)=0$ となる点が区間の内側に存在する.さらに,その区間で $h'>0$ (または $h'<0$)なら $h$ は狭義に単調だから,\,$0$ をとる点は高々1つである.存在することと高々1つであることを合わせて,解はただ1つと分かる.

不等式では,\,$h$ の最小値を調べる.増減表から最小値が $0$ 以上と分かれば,\,$h(x)\ge0$ を得る.等号が成り立つのは,\,$h$ が最小値 $0$ をとる点に限られる.

\begin{bookpoint}
\textbf{POINT}\par
差 $h=f-g$ を作る.方程式は端点での符号の違いで解の存在を,\,$h'$ の符号で個数を示す.不等式は $h$ の最小値が $0$ 以上であることを示す.
\end{bookpoint}

\begin{example}{方程式と不等式}
\begin{enumerate}[label=(\arabic*),leftmargin=*]
 \item 方程式 $x=\cos x$ が区間 $[0,1]$ にただ1つの解をもつことを示せ.
 \item $x>0$ のとき $\log x\le x-1$ を示し,等号成立条件を求めよ.
\end{enumerate}
\end{example}
\paragraphbf{方針}
(1) は $h(x)=x-\cos x$ と置き,中間値の定理で存在を,\,$h'$ の符号で個数を示す.(2) は $g(x)=x-1-\log x$ と置き,その最小値が $0$ であることを示す.

\paragraphbf{解法}
\begin{enumerate}[label=(\arabic*),leftmargin=*]
 \item $h(x)=x-\cos x$ と置く.\,$h$ は $[0,1]$ で連続であり,
 \[
  h(0)=-1<0,\qquad h(1)=1-\cos1>0
 \]
 である.中間値の定理より,\,$h(x)=0$ となる点が $(0,1)$ に存在する.一方,\,$0\le x\le1$ では
 \[
  h'(x)=1+\sin x>0
 \]
 だから,\,$h$ は $[0,1]$ で狭義に増加し,\,$0$ をとる点は高々1つである.したがって,方程式 $x=\cos x$ は $[0,1]$ にただ1つの解をもつ.
 \item $g(x)=x-1-\log x$ ($x>0$) と置く.
 \begin{align*}
  g'(x)&=1-\frac1x\\
  &=\frac{x-1}{x}
 \end{align*}
 であり,\,$x>0$ だから $g'$ の符号は $x-1$ の符号と一致する.よって $g$ は $0<x\le1$ で狭義に減少し,\,$x\ge1$ で狭義に増加する.したがって $g$ は $x=1$ で最小値 $g(1)=0$ をとり,\,$g(x)\ge0$,すなわち $\log x\le x-1$ を得る.\,$x\ne1$ では $g(x)>g(1)=0$ だから,等号は $x=1$ のときに限る.
\end{enumerate}

同じ方法で,後の節でも使う不等式をもう1つ示しておこう.\sectionref{節:等比級数}では,\,$c>0$ と自然数 $n$ に対する不等式 $(1+c)^n\ge1+nc$ を使い,そこで,指数が実数の場合は微分を学んだ後に示すと予告した.次の例題がその不等式である.

\begin{example}{実数べきの不等式}\label{例題:実数べきの不等式}
$\alpha$ を $\alpha\ge1$ を満たす実数とする.\,$x>-1$ のとき
\[
 (1+x)^\alpha\ge1+\alpha x
\]
が成り立つことを示せ.
\end{example}
\paragraphbf{方針}
左辺から右辺を引いた $h(x)=(1+x)^\alpha-1-\alpha x$ を作り,その最小値が $0$ であることを示す.\,$h'$ の符号は,\,$1+x$ が $1$ より大きいか小さいかで決まる.

\paragraphbf{解法}
$h(x)=(1+x)^\alpha-1-\alpha x$ ($x>-1$) と置く.\,$x>-1$ では $1+x>0$ だから,実数べきの微分の公式が使えて,
\begin{align*}
 h'(x)&=\alpha(1+x)^{\alpha-1}-\alpha\\
 &=\alpha\left\{(1+x)^{\alpha-1}-1\right\}
\end{align*}
となる.

まず $\alpha>1$ とする.このとき指数 $\alpha-1$ は正である.正の数を正の指数でべき乗した値は,底が $1$ より小さければ $1$ より小さく,底が $1$ より大きければ $1$ より大きい.したがって,
\begin{itemize}
 \item $-1<x<0$ では $0<1+x<1$ だから $(1+x)^{\alpha-1}<1$ であり,\,$h'(x)<0$ となる.
 \item $x>0$ では $1+x>1$ だから $(1+x)^{\alpha-1}>1$ であり,\,$h'(x)>0$ となる.
\end{itemize}
よって $h$ は $-1<x\le0$ で狭義に減少し,\,$x\ge0$ で狭義に増加するから,\,$x=0$ で最小値 $h(0)=1-1-0=0$ をとる.ゆえに $h(x)\ge0$,すなわち $(1+x)^\alpha\ge1+\alpha x$ である.等号は $x=0$ のときに限る.

$\alpha=1$ のときは,両辺がともに $1+x$ であり,つねに等号が成り立つ.

\medskip
この不等式で $\alpha$ を自然数 $n$,\,$x$ を正の数 $c$ とすれば,\,$(1+c)^n\ge1+nc$ となる.\,\sectionref{節:等比級数}の不等式は,この例題の特別な場合だったのである.指数を自然数から $1$ 以上の実数へ,\,$c>0$ を $x>-1$ へ広げても,同じ不等式が成り立つのだ.

では,指数が $1$ より小さい $0<\alpha<1$ の場合はどうなるだろうか.このときも $h'(x)=\alpha\left\{(1+x)^{\alpha-1}-1\right\}$ は変わらないが,指数 $\alpha-1$ が負になる.負の指数では大小が入れ替わり,\,$-1<x<0$ で $(1+x)^{\alpha-1}>1$,\,$x>0$ で $(1+x)^{\alpha-1}<1$ となる.つまり $h'$ の符号が $\alpha>1$ のときと逆になり,\,$h$ は $x=0$ で最大値 $0$ をとる.したがって,\,$0<\alpha<1$ では不等号が逆になって
\[
 (1+x)^\alpha\le1+\alpha x\qquad(x>-1)
\]
が成り立つ.
